\subsection{差速底盘与电机控制}\label{sec:chassis}

机器人靠两个直流电机驱动，一侧一个，并没有独立的转向舵机。这种结构称为差速驱动，前进
靠两侧轮速相同，转向则完全来自两侧轮速之差。两轮同速向前则直行，一侧快一侧慢则向轮速较
低的一侧偏转，两侧反向旋转则原地转身。感知给出的目标是球或桶在画面里的位置，控制要做的
就是把这个横向偏差换算成一对左右轮速，让机器人一边逼近目标一边把它拢到画面中央，最终停
在可以伸臂抓取的位置。

\repfig{motor-architecture.png}{从控制状态机到 H 桥的电机控制链路，左右轮速经电机抽象层映射为四路 PWM 占空比}{fig:chassis}

控制环的核心是一台状态机，它在追球、抓取、找桶、接近、投放几个状态之间流转，每一帧都输出
一份带着动作类型与左右轮速的指令。追球时它把球心相对画面中线的偏移作为转向偏置，偏移越大
偏置越强，远处采用两轮同向的差速行进，把偏置一加一减地叠加到一个随目标面积递减的基础速度
上，越近基础速度越低；当球已足够近却尚未对正抓爪时改为两轮反向的原地转身，使目标始终留在
视野内；一旦球进入设定的近距并落在抓爪正前方的窗口内连续若干帧，便切到制动并触发抓取。接
近桶的逻辑同出一辙，区别只在阈值与增益。这样状态机吐出的始终是一对范围在 $[-100, 100]$ 的有
符号轮速，正值代表该侧向前。

状态机之上是一层电机抽象，它把控制语义与具体硬件隔开。这层接口只有三个动作，\texttt{drive}
按左右两路设定轮速，\texttt{brake} 主动锁住两轮，\texttt{standby} 让输出断电滑行。控制环只与
这三个动作打交道，并在这一层做死区补偿，任何非零但过小的指令都会被抬到一个最小速度之上，否
则电机只发热而不转动。再往下，这一对有符号轮速被翻译成驱动 H 桥所需的占空比，电机随之在底盘
上实际转起来，机器人完成对球的追逐、抓取与投放。

\paragraph{从轮速到 H 桥的四路 PWM}
两个直流电机经一块 DRV8833 一类的 H 桥驱动板供电。H 桥用一对开关的通断决定电流方向，从而决定
电机正转、反转或停止，每个电机需要两条方向线，两个电机合起来就是四路带占空比调节的 PWM。占空
比的高低决定有效电压，也就决定了转速，而那对有符号轮速的符号则选定每个电机两条线中哪一条出
PWM、哪一条拉低，由此定下转向。一条轮速为正就让正转线按其大小出 PWM、反转线拉低，为负则两条
线互换，为零则两条线同时拉低让电机自由滑行；制动则把两条线同时拉高，让电机绕组短接而锁死。
四路占空比按这一规则从左右轮速算出，差速底盘的全部行为都落在这四个数上。

\paragraph{在 Starry 上提供与 Linux 一致的 PWM 通路}
\textbf{为了让上面这条链路在 Starry 与 Linux 上以同一份代码运行，Starry 提供了与 Linux 兼容的
\texttt{/sys/class/pwm} sysfs 接口，应用通过导出通道、写入周期与占空比、再启停通道来操控
每一路 PWM，无需关心底层寄存器。}在 RK3588 平台一侧，我们接好了对应的时钟、引脚复用与寄存器
配置，把 OrangePi 5 Plus 四十针排针上引出的 PWM 通道接到这套 sysfs 之下，于是同样的导出、
设周期与占空比、启停的操作序列在两个系统上行为一致，电机抽象层写出的四路占空比可以原样落到
硬件上。\textbf{该 PWM 驱动正在向上游提交。}

电机抽象与底层 PWM 之间是按值传递的轮速与占空比，而非跨越内核的复杂调用，因此从相机一帧到
落到电机上的那条指令始终走在纯计算与一次 sysfs 写入的轻量路径上，不会被驱动阻塞拖慢，与本
报告其余环节追求的帧到指令低时延保持一致。
